\begin{proposition}[隐藏位二层细分诱导的子群胚包含与指数界]\label{prop:op-algebra-hiddenbit-sheet-subgroupoid-index}
固定 \(m\ge 1\)，沿用引理 \ref{lem:pom-one-bit} 的隐藏位映射 \(b:\Omega_m\to\{0,1\}\)。
令 \(\mathcal R_m=\Omega_m\times_{X_m}\Omega_m\) 为折叠等价关系群胚（定理 \ref{thm:op-algebra-fold-basic-construction-pair-groupoid}），并定义其“同 sheet”子群胚
\[
\mathcal H_m
:=\{(\omega,\omega')\in\Omega_m\times\Omega_m:\ \Fold_m(\omega)=\Fold_m(\omega'),\ b(\omega)=b(\omega')\}
\ \subset\ \mathcal R_m.
\]
对每个 \(x\in X_m\) 记纤维 \(F_x:=\Fold_m^{-1}(x)\)，并令
\[
d_{m,i}(x):=\card{\{\omega\in F_x:\ b(\omega)=i\}}\qquad(i\in\{0,1\}),
\qquad d_m(x)=d_{m,0}(x)+d_{m,1}(x).
\]
则卷积 \(\ast\)-代数在半单分解下满足
\[
\CC[\mathcal R_m]\ \cong\ \bigoplus_{x\in X_m} M_{d_m(x)}(\CC),
\qquad
\CC[\mathcal H_m]\ \cong\ \bigoplus_{x\in X_m}\Bigl(M_{d_{m,0}(x)}(\CC)\oplus M_{d_{m,1}(x)}(\CC)\Bigr),
\]
且包含 \(\CC[\mathcal H_m]\subset \CC[\mathcal R_m]\) 在每个 \(x\)-块上对应于
\(
M_{d_{m,0}(x)}(\CC)\oplus M_{d_{m,1}(x)}(\CC)\hookrightarrow M_{d_m(x)}(\CC)
\)
的块对角嵌入。
因此若以向量空间维数比定义该包含的有限维指数
\[
[\CC[\mathcal R_m]:\CC[\mathcal H_m]]
:=\frac{\dim_\CC \CC[\mathcal R_m]}{\dim_\CC \CC[\mathcal H_m]}
=\frac{\sum_{x\in X_m}d_m(x)^2}{\sum_{x\in X_m}\bigl(d_{m,0}(x)^2+d_{m,1}(x)^2\bigr)},
\]
则必有双边界
\[
\boxed{\;
1\ \le\ [\CC[\mathcal R_m]:\CC[\mathcal H_m]]\ \le\ 2\; }.
\]
并且上界 \(2\) 取到当且仅当对所有 \(x\in X_m\) 都有 \(d_{m,0}(x)=d_{m,1}(x)\)（每条稳定类型纤维在两张 sheet 上严格等势分裂）。
\leanverified{paper\_op\_algebra\_hiddenbit\_sheet\_subgroupoid\_index}
\end{proposition}
\begin{proof}
半单分解与块嵌入由定理 \ref{thm:op-algebra-fold-basic-construction-pair-groupoid} 给出的纤维成对群胚实现逐块推出：
\(\mathcal H_m\) 在每条纤维 \(F_x\) 上分解为两份成对群胚
\(
F_{x,0}\times F_{x,0}\ \sqcup\ F_{x,1}\times F_{x,1}
\)，其中 \(F_{x,i}:=\{\omega\in F_x:\ b(\omega)=i\}\)。
于是对应代数为两块矩阵代数的直和，且包含为块对角嵌入。
\par
维数公式为 \(\dim_\CC M_d(\CC)=d^2\)，从而
\(
\dim \CC[\mathcal R_m]=\sum_x d_m(x)^2
\)、
\(
\dim \CC[\mathcal H_m]=\sum_x(d_{m,0}(x)^2+d_{m,1}(x)^2)
\)。
逐点不等式
\[
d_{m,0}(x)^2+d_{m,1}(x)^2\ \le\ (d_{m,0}(x)+d_{m,1}(x))^2\ \le\ 2\bigl(d_{m,0}(x)^2+d_{m,1}(x)^2\bigr)
\]
给出
\(
\sum_x(d_{m,0}(x)^2+d_{m,1}(x)^2)\le \sum_x d_m(x)^2\le 2\sum_x(d_{m,0}(x)^2+d_{m,1}(x)^2)
\)，从而得到 \(1\le[\cdot:\cdot]\le 2\)。
上界取到当且仅当上述右侧不等式对每个 \(x\) 逐项取等，即 \(d_{m,0}(x)=d_{m,1}(x)\)。
证毕。
\end{proof}

\endinput
